\subsection{An ordinary infinite-support pair first excluded at class three}
\label{sec:ghost-class3-example}

\begin{theorem}[First unrestricted integral failure at class three]
\label{thm:ghost-first-integral-class3}
There is a recursively specified ordinary pair
\[
 A=3\prod_{j\ge0}q_j^{5(j+4)!},\qquad
 B=-2\prod_{j\ge0}r_j^{5(j+4)!},
\]
with pairwise distinct primes \(q_j,r_j>3\), for which
the first empty fibre of global continuous integral
nilpotent sections is class three. Its class-two fibre
has exactly twenty-four geometric-conjugacy classes.
Every one has a nonzero local third obstruction at
\(q_0\) in its coefficient-prime-five component.
The reduction modulo \(25\) of that obstruction is
already nonzero.
\end{theorem}
\begin{proof}
We give the prime condition, prove that it is compatible
with the ordinary-pair recursion, and then calculate the
obstruction of every integral class-two lift.

\emph{A finite field extension attached to the point three.}
Let \(H\) be the exponent-five Heisenberg group with
coordinates \((a,b,h)\in\mathbb F_5^3\) and product
\[
 (a,b,h)(a',b',h')=(a+a',b+b',h+h'+ab').
\]
Its center and commutator subgroup are
\(Z=\{(0,0,h):h\in\mathbb F_5\}\).
Reducing the actual class-two geometric group and its
tangential arithmetic action gives
\[
 H\rtimes\mathbb F_5^\times,
 \qquad \lambda(a,b,h)=(\lambda a,\lambda b,\lambda^2h).
\]
The rational point \(3\) supplies a homomorphism
\(\rho_3:G_{\mathbb Q}\to H\rtimes\mathbb F_5^\times\).
Write \(L\) for the fixed field of its kernel and put
\(K=\mathbb Q(\mu_5)\). Its restriction to \(G_K\)
is onto \(H\). Indeed its abelianization is the pair
of Kummer characters of \(3,-2\), independent over
\(K\): valuations at primes above three and two
separately detect the two exponents modulo five.
These primes are unramified in \(K/\mathbb Q\).
A subgroup of \(H\) mapping onto \(\mathbb F_5^2\)
contains two elements with independent abelianizations;
their commutator is a nonzero element of \(Z\).
It therefore contains all of \(Z\) and is all of \(H\).
Since the cyclotomic character is onto,
\[
 \operatorname{Gal}(L/\mathbb Q)=H\rtimes\mathbb F_5^\times.
\]
The commutator subgroup of this group is exactly \(H\).
It is contained in \(H\), and its abelian image contains
\((\lambda-1)\mathbb F_5^2\) for \(\lambda=2\),
hence all of \(\mathbb F_5^2\); the preceding commutator
argument then gives all of \(H\). Consequently
\[
                  L^{\rm ab}=K.                 \tag{1}
\]
The field \(L\) and the condition that a Frobenius be
a nonzero member of \(Z\) do not depend on a choice
of geometric path to the point: changing the path
conjugates \(\rho_3\), leaving its kernel and this
conjugacy-invariant condition unchanged.

Here is an explicit radical presentation of this particular
field, making the prime condition below effective. Fix
\(\zeta=\zeta_5\), and put
\[
 \begin{gathered}
 X^5=3,\qquad Y^5=-2,\qquad
 \Theta(X)=\prod_{i=0}^4(1-\zeta^iX)^i,\\
 Z_0^5=\Theta(X),\qquad
                 L=\mathbb Q(\zeta,X,Y,Z_0).
 \end{gathered}                                      \tag{1a}
\]
We prove that this is the field defined by \(\rho_3\),
not just an abstract Heisenberg extension with the same
abelian quotient. First replace three by an indeterminate
\(t\), and minus two by \(1-t\). Over \(K(t)\)
the radical cover has automorphisms
\[
 \begin{array}{c|ccc}
       &X&Y&Z_0\\\hline
 \sigma&\zeta X&Y&Z_0(1-X)/Y\\
 \tau&X&\zeta Y&Z_0.
 \end{array}                                        \tag{1b}
\]
Their defining identity is
\[
 \frac{\Theta(\zeta X)}{\Theta(X)}
       =\frac{(1-X)^5}{1-X^5}.
\]
The class of \(\Theta(X)\) is not a fifth power in
\(K(X,Y)\). Otherwise the commuting automorphisms
\(\sigma,\tau\) of that abelian extension of \(K(t)\)
would act on a fifth root as
\(\sigma(Z_0)=\zeta^a Z_0(1-X)/Y\) and
\(\tau(Z_0)=\zeta^b Z_0\), for some integers
\(a,b\). These formulas give
\(\sigma\tau(Z_0)=\zeta\tau\sigma(Z_0)\), a
contradiction. Thus the cover has degree \(125\).
The displayed lifts satisfy \(\sigma^5=\tau^5=1\)
and \([\sigma,\tau](Z_0)=\zeta Z_0\), with that
commutator central. Its group is precisely \(H\).

The divisors of the radical expressions are supported
over \(t=0,1,\infty\), so the cover is finite etale
over \(U_K\). It has the specified tangential branch
at zero
\[
           X=t^{1/5},\qquad Y(0)=Z_0(0)=1.
\]
Local expansions at zero give the generator \(\sigma\)
of the inertia above zero. On the branch \(X\to1\)
at one, \(\Theta(1)\ne0\), because its zero-exponent
factor is omitted; the inertia generator there is
\(\tau\). Choosing the intervening geometric path
only changes the resulting marking by the usual inner
ambiguity. The free geometric group on these two
inertias has a class-two exponent-five quotient of
order \(5^3\): every element collects uniquely to
\(\tau^b\sigma^a[\sigma,\tau]^h\). The surjection
from this quotient to the degree-\(125\) cover is
therefore an isomorphism. Thus this cover is the
actual geometric quotient used to define \(\rho_3\),
with its tangential trivialization.

The descent and arithmetic action can also be written
explicitly. For \(\lambda\in\mathbb F_5^\times\)
take \(k\in\{1,2,3,4\}\) representing
\(\lambda^{-1}\). For \(j=0,\ldots,4\), let
\(i_j\in\{0,\ldots,4\}\) be the representative
of \(\lambda^{-1}j\). Define
\[
 P_\lambda(X)=
       \prod_{j=0}^4(1-\zeta^jX)^{(i_j-kj)/5}.
\]
Every exponent is an integer, and reindexing gives
\[
 \Theta(X)\big|_{\zeta\mapsto\zeta^\lambda}
             =\Theta(X)^kP_\lambda(X)^5.
\]
Hence \(\zeta\mapsto\zeta^\lambda\), \(X,Y\mapsto X,Y\),
\(Z_0\mapsto Z_0^kP_\lambda(X)\) defines a semilinear
automorphism of the generic cover. These automorphisms
compose exactly: their possible discrepancy fixes
\(X,Y,\zeta\), hence multiplies \(Z_0\) by a fifth
root of unity; preservation of its tangential constant
term one makes that root equal to one. Conjugation sends
\(\tau\) to \(\tau^\lambda\) directly. It sends
\(\sigma\) to \(\sigma^\lambda\): their actions on
\(X,Y\) agree, and the same constant-term argument
removes the possible central factor. The action on
\([\sigma,\tau]\) is therefore multiplication by
\(\lambda^2\), exactly the arithmetic action above.
Specialization at \(t=3\) is the pullback along the
point-three section, so gives (1a). The specialized
Kummer classes remain independent by the valuations
already checked; consequently its image is still the
full group \(H\rtimes\mathbb F_5^\times\).
The exact radical identities and this finite group's
order and commutator subgroup are independently checked
in \path{code/check_heisenberg_point3.py}.

\emph{The compatible first prime.}
Use the A.44 recursion with
\[
 n_j=5(j+4)!,\quad a_0=3,\quad b_0=-2,
 \quad H_j(p)=2v_p(n_{j+1})+j+3,
 \quad B(n)=(n+1)^4.
\]
In its first stage \(n_0=120\), \(n_1=600\).
Choose every old local solution to be \(X_p=Y_p=1\),
which is an exact solution of \(3X_p^{120}-2Y_p^{120}=1\).
Thus the first prime's CRT conditions are all
\(q_0\equiv1\pmod{p^{H_1(p)}}\). Let \(M\)
be their product over the first old set, namely the
primes at most \(B(600)\) together with the divisors
of \(a_0b_0\). Put
\[
 D=\operatorname{lcm}(M,120,25),\qquad
 F=\mathbb Q(\mu_D),\qquad
 R=F\bigl((-2)^{1/120},3^{1/5}\bigr).
\]
This is a finite Galois extension of \(\mathbb Q\),
and \(R/F\) is abelian, since \(F\) contains
\(\mu_{120}\). Equation (1) gives
\(L\cap F=K\); thus
\(\operatorname{Gal}(LF/F)=H\).
The intersection \(LF\cap R\) is an abelian
extension of \(F\), as a subfield of \(R\).
It is therefore fixed by the commutator subgroup
\(Z\) of \(H\). A chosen nonzero \(z\in Z\)
and the identity of \(R\) agree on this intersection,
so extend to an automorphism of \(LR\) fixing \(R\).
Chebotarev's theorem
\cite[Theorem 8.31]{GhostMilneANT} supplies infinitely
many unramified rational primes in its conjugacy class.
For each of them
\[
 \begin{gathered}
 q\equiv1\pmod D,\qquad q\text{ splits completely in }R,\\
 \rho_3(\operatorname{Frob}_q)=(0,0,c;1),
                         \qquad c\in\mathbb F_5^\times.
 \end{gathered}                                      \tag{2}
\]
Conjugation can replace \(c\) by \(\lambda^2c\),
which is still nonzero. Choose \(q_0\) to be the
least such prime exceeding all the first-stage A.44
bounds and outside the finite ramification set.
This specifies a finite prime-selection condition
without claiming a numerical value for this very large
prime. In particular \(q_0\equiv1\pmod{25}\).

Equivalently, this is an entirely finite arithmetic
prime test. Besides the lower bound, test
\[
 q\equiv1\pmod D,\qquad
 (-2)^{(q-1)/120}=1\pmod q,\qquad
 3^{(q-1)/5}=1\pmod q.
\]
Choose an element \(\zeta_q\in\mathbb F_q\) of
order five and a root \(x_q^5=3\), and form
\[
 \theta_q=\prod_{i=0}^4(1-\zeta_q^ix_q)^i.
\]
The remaining condition is exactly
\[
                   \theta_q^{(q-1)/5}\ne1.          \tag{2a}
\]
Indeed the first three tests say that the prime splits
in \(R\); over this split base, the last radical in
(1a) has nontrivial central Frobenius precisely when
\(\theta_q\) is not a fifth power. Replacing
\(x_q\) by \(\zeta_qx_q\) multiplies \(\theta_q\)
by \((1-x_q)^5/(-2)\), a fifth power because
\(-2\) is already a fifth power at this prime.
Replacing \(\zeta_q\) by another generator raises
its fifth-power class to a unit power, by the explicit
descent identity. Thus (2a) is independent of these
choices. Trial primality, powering, and root extraction
make the least-prime rule algorithmic.
The retained certificate
\path{code/check_heisenberg_point3_prime.py}
checks the same radical and central-Frobenius tests
with the smaller modulus \(600\). Its prime is a
finite illustration of these tests, not the recursively
defined \(q_0\), whose old-prime CRT modulus is much
larger.

These requirements imply every original first-prime
condition: the old CRT classes are one, and the prime
splits in \(\mathbb Q(\mu_{120},(-2)^{1/120})\).
The next prime \(r_0\), and all subsequent pairs of
primes, are chosen by the original A.44 rules, taking
the least eligible primes and the least eligible finite
residue choices to fix the recursion.

For precision, replacing \((j+4)!\) by \(5(j+4)!\)
preserves all inputs of that proof. Each exponent is
divisible by twenty-four, the ratio
\(n_{j+1}/n_j=j+5\) is unchanged, and both radicands
retain the forms \(3w^8\) and \(-2w^8\), with
\(\gcd(w,6)=1\). The radical-field intersection
calculation therefore applies unchanged at each later
prime choice. At the old primes, the same bound
\(H_j(p)=2v_p(n_{j+1})+j+3\) gives the Hensel lifts
and the congruences needed for the next exponent.
At all remaining primes the same finite-field estimate
with \(B(n)=(n+1)^4\) gives a unit solution.
The resulting sequences
\[
 a_{j+1}=a_jq_j^{n_j},\qquad
 b_{j+1}=b_jr_j^{n_j}
\]
are Cauchy in every global power quotient, since
\(n_j\) is eventually divisible by every positive
integer. At a fixed finite prime, their valuations
eventually stabilize, their unit parts converge, and
\(v_p(a_j+b_j-1)\) tends to infinity. Thus their
limits are ordinary nonzero local numbers summing
to one, with valuation vectors
\[
 \begin{array}{c|cc}
  p&v_p(A)&v_p(B)\\\hline
  3&1&0\\
  2&0&1\\
  q_j&n_j&0\\
  r_j&0&n_j\\
  \text{otherwise}&0&0.
 \end{array}
\]
This verifies ordinarity and infinite support for the
modified recursion. The extra condition (2) affected
only a compatible first-prime choice.

\emph{The reduction of every integral class-two lift.}
Let \(s\) be any such lift and restrict to coefficient
prime five. A geometric conjugation makes its two
abelian cocycles \(a,b\) reduce literally to those
of the point-three section modulo five: their Kummer
classes agree, since every new exponent is divisible
by five. A lift of the required abelian conjugator
exists in the class-two group. The difference between
its central coordinate \(h\) and the point-three
central coordinate \(h_3\), modulo five, is now a
literal cocycle
\[
            w=h-h_3\in Z^1(G_{\mathbb Q},\mathbb F_5(2)).
\]
It is unramified outside five. We verify this for the
actual integral restrictions, before reducing them.
At a new prime with valuation vector \((m,0)\),
the ordinary complement is a principal unit. Thus
\(b=\delta u\) for an integral local
\(u\in\mathbb Z_5(1)\). With a tame Frobenius
\(F_q\) and inertia generator \(T_q\), write
\(\beta=b(F_q)=(q-1)u\). The class-two tame
relation \(F_qT_qF_q^{-1}=T_q^q\), using the
displayed Heisenberg law integrally, gives
\[
 q(q-1)h(T_q)=qm\beta,
              \qquad h(T_q)=mu.                 \tag{3}
\]
It is zero modulo five. At a new prime with valuation
\((0,m)\), the same calculation with
\(a=\delta v\) gives \(h(T_q)=-mv\), again
zero modulo five. At a prime of zero valuation,
\((q^2-q)h(T_q)=0\) in the torsion-free module
\(\mathbb Z_5(2)\), so \(h(T_q)=0\).
The point-three integral section has zero inertia
at all these primes as well. At two and three the
difference \(w\), already a weight-two cocycle,
satisfies \((q^2-q)w(T_q)=0\). Both scalars are
units modulo five, hence its inertia is zero there.
Wild and non-five tame inertia vanish because the
integral geometric inertia is pro-five and the
arithmetic action is unramified. This proves the
claimed unramifiedness of \(w\).
Since the coefficient action is unramified outside five,
vanishing on these inertia groups also gives vanishing on
their closed normal closure. Thus \(w\) factors through
the actual quotient \(G_{\mathbb Q,\{5,\infty\}}\).

In fact
\[
 H^1(G_{\mathbb Q,\{5,\infty\}},\mathbb F_5(2))=0.       \tag{4}
\]
Inflation injects integral first cohomology into
\(H^1(G_{\mathbb Q},\mathbb Z_5(2))=0\), by
Proposition~\ref{prop:ghost-h1-weight2}. Also
\(H^2(G_{\mathbb Q,\{5,\infty\}},\mathbb Z_5(2))
=K_2(\mathbb Z)\{5\}=0\), by
\cite[VI.8.2, Table VI.10.1.1]{GhostWeibelKbook},
as used in Proposition~\ref{prop:ghost-ordinary-cup}.
The coefficient long exact sequence for multiplication
by five gives (4). Therefore \(w\) is a coboundary.
At \(q=q_0\), where \(q\equiv1\pmod5\), its
value on Frobenius is zero. Equation (2) now gives
\[
 \alpha=a(F_q)\equiv0\pmod5,\qquad
 \gamma=h(F_q)\equiv c\not\equiv0\pmod5.           \tag{5}
\]

\emph{The nonzero third obstruction.}
At \(q=q_0\), put \(m=120\). Its ordinary
valuation vector is \((m,0)\), so the notation
of Proposition~\ref{prop:ghost-class3-residue} has
\[
 n=0,\quad \beta=(q-1)u,\quad t=mu.
\]
Substitution in its first residue gives the exact
integral identity
\[
 Z_{D,q}
   =m(\gamma-q\alpha u)
       +(q^2-1)u\frac{m(m-1)}2.                 \tag{6}
\]
By (5) the first parenthesis is a five-adic unit,
and \(v_5(m)=1\). The second term has valuation
at least two, since \(q\equiv1\pmod{25}\).
Thus \(v_5(Z_{D,q})=1\), whereas
\(v_5(q^2-1)\ge2\). Its class in
\[
 H^2(G_{\mathbb Q_q},\mathbb Z_5(3))
             =\mathbb Z_5/(q^2-1)\mathbb Z_5
\]
is nonzero. Equivalently, the actual tame cochain map
of Proposition~\ref{prop:ghost-tame-cochain} gives
a nonzero residue; the same calculation modulo
twenty-five remains nonzero. Hence this specified
integral class-two section has no class-three lift.
The argument applied to an arbitrary \(s\), so it
excludes every integral class-two section above the
constructed pair. Class-two existence and its exact
twenty-four conjugacy classes follow from
Corollaries~\ref{cor:ghost-class2-lift} and
\ref{cor:ghost-class2-fibre24}. Class one is already
the pair itself. This proves the claimed first depth.
\end{proof}

The modulo-twenty-five assertion has the precise domain
used in the proof: it detects the third obstruction
of the reduction of each \emph{integral} class-two
section. The construction has not classified arbitrary
finite-coefficient class-two sections. Those admit
additional classes in
\(H^1(G_{\mathbb Q},\mathbb Z/25(2))\), and their
ramification need not satisfy (3) integrally. The
theorem therefore makes no claim that one finite
geometric quotient excludes all finite-coefficient
sections of this pair.
